\documentclass{report}
\usepackage{amsmath,amssymb}
\setlength{\parindent}{0mm}
\setlength{\parskip}{1em}
\begin{document}
\begin{center}
\rule{6in}{1pt} \
{\large Hasan Metin Aktulga \\
{\bf Solving Large-scale Eigenvalue Problems in Nuclear Configuration Interaction Calculations}}

Lawrence Berkeley National Lab \\ 1 Cyclotron Rd \\ MS 50F-1650 \\ Berkeley \\ CA 94720
\\
{\tt hmaktulga@lbl.gov}\\
Chao Yang\\
Esmond G. Ng\\
Pieter Maris\\
James P. Vary\end{center}

The structures of nuclei can be analyzed by computing the eigenvalues and
eigenvectors of a Hamiltonian matrix $H$ constructed by what is known in
quantum mechanics as the configuration interaction (CI) technique.In most
cases, we are interested in the lowest few eigenvalues of $H$. In some
applications, we may be interested in a relatively large number of low
energy states (eigenvalues) with a prescribed total angular momentum,
$J$. This type of calculation, which we will refer to as a total-J
calculation, is particularly useful for investigating, among others,
nuclear level densities as a function of $J$ and excitation energy and
for evaluating scattering amplitudes as a function of channel $J$. To
achieve this goal, we perform a simultaneous diagonalization of $H$ and
$J^2$, where $J^2$ is the total angular momentum square operator.

The Hamiltonian and the angular momentum matrices constructed from the
configuration interaction approximation are sparse. However, their
dimensions can be extremely large. Therefore, the matrix eigenvalue
problem must be solved by iterative methods.

In this talk, we discuss how to efficiently solve this type of
eigenvalue problem on multi-core machines with tens of thousands
of processing units. In particular, we will examine efficient
strategies for distributing matrices to achieve good load balance,
the communication overhead associated with orthogonalization and
efficient ways to perform sparse matrix vector multiplications.

In addition to discussing these issues in the context of the Lanczos
algorithm which is currently used in several nuclear CI codes, we will
also examine how we may benefit from the use of a block iterative solver
such as a polynomial accelerated subspace
iteration or the locally optimal block conjugate gradient (LOBPCG)
method, especially when the Hamiltonian is so large that it cannot
be completely stored in the main memory.


\end{document}
